\chapter{Risultati}
\label{cha:results}

In questo capitolo sono presentati e discussi i risultati sperimentali emersi dall'esecuzione della pipeline descritta nel Capitolo~\ref{cha:methodology}. In particolare vengono indagate le \textit{research questions}:
\begin{enumerate}[label=\textbf{RQ\arabic*}]
    \item La codifica di features umane influenza la qualità della predizione sul risultato di montaggio?
    \item L'uso di una rete neurale ricorrente supera le performance dei modelli ad albero?
    \item Quali features motivano le predizioni e come varia la loro importanza durante l'esecuzione del processo?
    \item Le tecniche di spiegabilità \textit{model-specific} sono sufficientemente fedeli a quelle agnostiche?
\end{enumerate}

\section{RQ1:}
\label{sec:rq1_tree_models}
Lo scopo della prima domanda di ricerca è determinare se l'arricchimento dei process log tradizionali con informazioni ergonomiche e posturali riguardanti l'operatore produce un aumento della capacità predittiva dei modelli tradizionalmente utilizzati per l'analisi di dati tabulari.

\subsection{Confronto della \textit{baseline} con metodi \textit{ensemble}}
\label{subsec:rq1_baseline_vs_ensemble}
Come descritto nel Capitolo~\ref{cha:methodology}, l'indagine sperimentale replica il lavoro di riferimento svolto da Vian~\cite{ppm-human-centered-vian2024}, in cui viene addestrato un singolo \textit{decision tree}.

Per superarne le limitazioni descritte nel Capitolo~\ref{cha:state_of_art} sono stati impiegati i modelli \textit{ensemble}
\begin{enumerate*}[label=(\roman*)]
    \item Random Forest (che applica la tecnica del \textit{bagging}) e
    \item XGBoost (che applica la tecnica del \textit{boosting})
\end{enumerate*}, mantenendo il Decision Tree come riferimento per l'analisi statistica.


\subsection{Confronto configurazioni \texttt{no\_human} vs \texttt{human}}
\label{subsec:rq1_human_vs_nohuman}

\subsection{Miglior classificatore ad albero}
\label{subsec:rq1_top_classifier}

\subsection{Migliore strategia di \textit{encoding}}
\label{subsec:rq1_top_classifier}

\section{RQ2:}
\label{sec:rq2_rf_vs_lstm}

\subsection{Valutazione comparativa Random Forest vs LSTM}
\label{subsec:rq2_rf_vs_lstm_encoding}

\subsection{Trade-off tra complessità architetturale e performance}
\label{subsec:rq2_tradeoffs}


\section{RQ3:}
\label{sec:rq3_shap_explainability}

\subsection{Quota d'importanza delle feature di base}
\label{subsec:rq3_base_feature}

\subsection{Quota d'importanza temporale}
\label{subsec:rq3_step_dynamics}

\subsection{Importanza del distretto di feature RULA}
\label{subsec:rq3_rula_wrist}

\section{R4:}
\label{sec:rq4_methodological_comparison}

\subsection{Fedeltà delle spiegazioni}
\label{subsec:rq4_fidelity_convergence}

\subsection{Trade-off di efficienza computazionale}
\label{subsec:rq4_computational_efficiency}